\documentclass[12pt,a4paper]{article}
\usepackage[margin=2cm]{geometry}
\usepackage{xeCJK}
\usepackage{fontspec}
\setCJKmainfont{Noto Serif CJK TC}[Script=CJK]
\usepackage{amsmath,amssymb}
\usepackage{graphicx}
\usepackage{fancyhdr}
\setlength{\headheight}{14.5pt}
\addtolength{\topmargin}{-2.5pt}
\usepackage{hyperref}
\usepackage{listings}
\usepackage{enumitem}
\usepackage{titlesec}
\usepackage{caption}
\usepackage{indentfirst}
\usepackage{booktabs}
\usepackage{longtable}
\usepackage{multirow}
\usepackage{array}
\usepackage{tabularx}
\usepackage{float}
\usepackage{minted}
\usepackage{titling}
\newcommand{\subtitle}[1]{%
  \posttitle{%
    \par\end{center}
    \begin{center}\large #1\end{center}
    \vskip 0.5em}%
}
\setlength{\parindent}{2em}
\pagestyle{fancy}
\fancyhf{}
\cfoot{\thepage}
\linespread{1.3}
\setminted{
    linenos,                % 行號
    frame=lines,            % 上下框線
    framesep=5pt,           % 程式碼與邊框距離
    numbersep=8pt,          % 行號與程式碼距離
    fontsize=\scriptsize,   % 字體大小
    breaklines,             % 自動換行
    tabsize=4,              % tab 寬度
    rulecolor=\color{black},% 框線顏色
    xleftmargin=1.5em       % 左側縮排
}


\title{Is Eastern Happiness All the Same? Comparing Conceptions of Happiness in Taiwan, Japan, and Canada}
\subtitle{Well-Being And Life Final Report Proposal}
\author{B12508026戴偉璿}
\date{\today}

\begin{document}

\maketitle

\lhead{Well-Being And Life Final Report Proposal}
\rhead{B12508026戴偉璿}
\newpage
\tableofcontents
\newpage


\begin{abstract}
It is commonly believed that students in East Asia are less happy than their Western counterparts, as East Asian countries tend to rank lower in international happiness surveys. This report examines whether this belief holds, or whether the gap reflects different cultural conceptions of happiness rather than a genuine difference in well-being. Through a literature review, I compare Taiwan and Japan within East Asia, and contrast both with Canada as a Western case. Interviews with friends living in these countries are used to illustrate how the findings appear in everyday life.
\end{abstract}

\section{Introduction}
\emph{
A means average, B means below average, C means can't eat dinner, D means don't go home, F means find a new family.
}

Above is a common joke on the Internet to ridicule the rigid standard for Asian students.Though it is only a joke, but it also implies the Asian students' stress may be harder than students in other areas. High stress may decrease the happiness. In this research, I want to 

\section{Methods}

\end{document}